% TOP_PROBLEM: {"id": 30006686, "title": "Basing One-Way Functions on NP-Hardness", "category_id": 15, "category": {"id": 15, "name": "computer_science", "display_name": "Computer Science", "description": "Computational complexity, algorithms, and theoretical CS.", "slug": "computer-science", "order_index": 15, "created_at": "2026-07-31T15:26:25.671Z"}, "status": "open"}
% FIRST_POSTED: 2026-09-27
% !TeX program = lualatex
% Standalone notebook; compile twice: lualatex 30006686.tex
% Original section preserved verbatim except for marked insertion blocks.
\documentclass[11pt,a4paper]{article}
\usepackage[margin=24mm,headheight=14pt]{geometry}
\usepackage{fontspec}
\setmainfont{Latin Modern Roman}
\setsansfont{Latin Modern Sans}
\setmonofont{Latin Modern Mono}
\usepackage{amsmath,amssymb,mathtools}
\usepackage{unicode-math}
\setmathfont{Latin Modern Math}
\usepackage{microtype}
\usepackage{enumitem,xurl,needspace}
\usepackage[dvipsnames]{xcolor}
\usepackage{hyperref}
\hypersetup{unicode=true,colorlinks=true,linkcolor=MidnightBlue,urlcolor=MidnightBlue,
 pdftitle={Research notebook: Problem 30006686},
 pdfauthor={Open Problems Math research notebook},bookmarksnumbered=true}
\usepackage{bookmark,titlesec,fancyhdr}
\titleformat{\section}{\Large\bfseries\raggedright}{\thesection}{0.7em}{}
\pagestyle{fancy}\fancyhf{}
\fancyhead[L]{\small\sffamily Open Problems Math\enspace|\enspace Research notebook}
\fancyhead[R]{\small\sffamily Problem 30006686}
\fancyfoot[L]{\small\sffamily 27 September 2026}
\fancyfoot[R]{\small\sffamily \thepage}
\setlength{\parindent}{0pt}
\setlength{\parskip}{5pt plus 1pt}
\setlength{\emergencystretch}{3em}
\setcounter{secnumdepth}{1}
\setcounter{tocdepth}{1}
\setlist[itemize]{leftmargin=3em,itemsep=5pt,topsep=4pt,parsep=0pt}
\allowdisplaybreaks
\newcommand{\N}{\mathbb{N}}
\newcommand{\Z}{\mathbb{Z}}
\newcommand{\Q}{\mathbb{Q}}
\newcommand{\R}{\mathbb{R}}
\newcommand{\C}{\mathbb{C}}
\newcommand{\F}{\mathbb{F}}
\newcommand{\A}{\mathbb{A}}
\newcommand{\PP}{\mathbb P}
\newcommand{\E}{\mathbb{E}}
\newcommand{\eps}{\varepsilon}
\newcommand{\dd}{\,\mathrm d}
\newcommand{\sref}[2]{\hyperlink{source:#1:#2}{[S#2]}}
\newcommand{\eref}[2]{\hyperlink{extra:#1:#2}{[E#2]}}
\newif\ifshowcatalogue
\showcataloguetrue
\newcommand{\cataloguescope}[1]{\ifshowcatalogue
 \paragraph{Exact scope in the supplied catalogue.}
 \begin{quote}\small #1\end{quote}\fi}
\numberwithin{equation}{section}
\begin{document}
\setcounter{section}{0}
% BEGIN PRESERVED SECTION
% ================================================================
% problemId: 30006686
% Canonical problem ID: 30006686
% exactTarget SHA-256: 80f98f05400696493a54220aee8bf098a70484b009e4343fe292d14aafa57319
\clearpage
\section*{\texorpdfstring{Basing One-Way Functions on NP-Hardness}{Basing One-Way Functions on NP-Hardness}}\label{30006686}
\subsection{Definitions and mathematical statement}
A worst-case hardness assertion excludes an efficient algorithm correct on every input; average-case inversion hardness bounds success on $f(U_n)$, where $U_n$ is uniform on $\{0,1\}^n$. The desired construction and reduction would establish
\[
 \mathsf P\ne\mathsf{NP}\quad\Longrightarrow\quad
 \exists\text{ a one-way function family}.
\]
More operationally, an efficient inverter with nonnegligible success would be transformed into an efficient solver for an NP-hard decision problem. The class of allowed reductions matters: an obstruction to a restricted black-box reduction does not by itself refute the unrestricted goal. The catalogue leaves this reduction model broad.
\par\smallskip\textit{Formulation and concept references:} \sref{30006686}{1}.
\cataloguescope{Determine whether one-way functions can be based on NP-hardness, i.e., whether there is a reduction from a worst-case NP-complete (or NP-hard) decision problem to the task of average-case inverting a polynomial-time computable function, so that P != NP would imply the existence of one-way functions.}
\subsection{Short English statement}
Can worst-case NP-hardness alone be turned into the average-case inversion hardness needed for one-way functions?
% BEGIN RESEARCH 33
\subsection{Research attempt}

\paragraph{Current frontier.}
Akavia--Goldreich--Goldwasser--Moshkovitz's nonadaptive-reduction obstruction
has important model hypotheses; it is not a disproof of the unrestricted
implication in this entry. Their adaptive result required a 2010 erratum.
Bogdanov--Brzuska subsequently prove, in Theorem~1, that a language admitting
their reduction to an approximately size-verifiable function belongs to
$\mathsf{AM}\cap\mathsf{coAM}$. The theorem and the unaffected nonadaptive
result, rather than an unqualified reading of the original abstract, are the
relevant barriers. \sref{30006686}{1}, \eref{30006686}{1}, \eref{30006686}{2}.

Hirahara's meta-complexity characterization assumes the NP-hardness of a
specific approximation problem in addition to worst-case NP hardness;
Liu--Pass's boundary characterization retains a substantial derandomization
assumption for ordinary time-bounded Kolmogorov complexity.
\sref{30006686}{3}, \sref{30006686}{4}. The 2026 zero-knowledge implication assumes
$\mathsf{NP}\not\subseteq\mathsf{ioP/poly}$ and existence of specified
nontrivial zero-knowledge protocols. It does not replace those hypotheses
by $\mathsf P\ne\mathsf{NP}$. \sref{30006686}{2}, \eref{30006686}{3}.
These are statement-level comparisons, not independent audits of the
complete recent proofs.

\paragraph{Attack route.}
I try to start from a transparent worst-case NP-hard inversion task and
amplify it by independent products. This separates the two jobs often
conflated in a proposed proof: encoding NP witnesses into an image fiber,
and ensuring that uniformly sampled inputs encounter hard fibers often
enough. The attempted construction below has the first property exactly.
I then prove that \emph{every} number of independent repetitions of this
particular construction still permits an efficient average-case inverter.
A second check isolates the difference between deterministic worst-case
hardness and resistance to randomized inversion.

\paragraph{Attempt.}
Fix a standard padded Boolean formula encoding: a description $z$ of $n$
bits has at most $n$ variables, and $V(z,w)$, for $w\in\{0,1\}^n$, evaluates
that assignment, ignoring unused bits. Invalid descriptions have value zero.
This is polynomial time, and every formula can be padded to such a format
with polynomial overhead. Define
\begin{equation}\label{r33:rare}
 G_{3n}(r,z,w)=
 \begin{cases}
  1\,\|\,r\,\|\,z\,\|\,w,&r\ne0^n,\\
  0\,\|\,0^n\,\|\,z\,\|\,V(z,w)\,\|\,0^{n-1},&r=0^n,
 \end{cases}
\end{equation}
where all three input blocks have $n\ge1$ bits. The output always has
$3n+1$ bits. At positive input lengths not divisible by three, use
$G_N(x)=1\|x$. Thus the evaluation algorithm is one uniform polynomial-time
algorithm on every length; inversion is allowed to return any preimage.

\textbf{Lemma 33.1 (hard worst-case fibers, easy uniform inversion).}
A polynomial-time algorithm that inverts $G$ correctly on every image point
would decide SAT in polynomial time. Nevertheless a deterministic inverter
succeeds on $G_{3n}(U_{3n})$ with probability at least $1-2^{-n}$.

\emph{Proof.} For a formula encoded by $z$, query the hypothetical inverter
on $0\|0^n\|z\|1\|0^{n-1}$. If the formula is satisfiable, this is an image
point and every valid preimage contains a satisfying $w$. If it is
unsatisfiable, no answer passes the direct check $V(z,w)=1$. One runs the
fixed polynomial-time inverter and verifies its answer, so no correctness
assumption on nonimage points is needed. This is an ordinary worst-case
reduction for the inversion task, not an average-case reduction.

For the average-case algorithm, an output starting in one displays its
entire original input. Return that input and fail on outputs starting in
zero. The probability of the latter branch is exactly $2^{-n}$, because
$r$ is uniform. This proves the claimed success bound.\hfill$\square$

Consequently ``inversion is NP-hard in the worst case'' does not imply that
\emph{the same function} is one-way. This example does not refute the
existence of another reduction or a different one-way function under
$\mathsf P\ne\mathsf{NP}$.

\medskip
\textbf{Lemma 33.2 (the independent-product repair also fails).}
Consider the componentwise product of $k\ge1$ copies of $G_{3n}$, with
$n,k$ given as parsing parameters. There is one inverter, polynomial in the
total input/output length $\Theta(nk)$, whose success on uniformly sampled
product inputs is bounded below by an absolute positive constant for every
$n,k$, and tends to one in the small-$k$ regime as $n\to\infty$.

\emph{Proof.} Compare the integers $k^2$ and $2^n$, a polynomial-time
operation on their binary representations.
If $k^2\le2^n$, invert each explicit branch as above and fail if any block
is on the rare branch. By the union bound,
\begin{equation}\label{r33:smallproduct}
 \Pr[\text{success}]=(1-2^{-n})^k
           \ge1-k2^{-n}\ge1-2^{-n/2}\ge1-2^{-1/2}>0.
\end{equation}
If $k^2>2^n$, enumerate the $2^n$ assignments for each rare block, using
$V$ to find a preimage with the indicated output bit. On sampled image
points such a preimage exists, including when that bit is zero. The work
is at most $k2^n\operatorname{poly}(n)<k^3\operatorname{poly}(n)$,
which is polynomial in $nk$. This branch succeeds with probability one.
Malformed inputs may be rejected immediately, and the enumeration stops
after its explicitly bounded list even on nonimage points. Hence the
algorithm has a worst-case polynomial running-time bound.\hfill$\square$

The lemma holds even when $k$ is exponentially large as a function of $n$:
in that regime brute-force search is affordable in the \emph{new} total
input length. It is not legitimate to call the brute force exponential
without specifying which security parameter is being charged. The result
is specific to this rare-core encoding and does not disprove general
hardness amplification, whose starting point needs an appropriate
inverse-polynomial hardness fraction.

\medskip
\textbf{Can a reduction ignore distributional exposure?}
No. If an inverter is perfect outside a set of image points of measure
$\varepsilon$ under $G(U)$, its behavior on that set can be arbitrary.
A reduction placing its important query inside that set cannot use the
average success guarantee alone. Even when every query marginal has the
same distribution as $G(U)$, nonnegligible inverter success does not imply
that a polynomial collection of queries contains all the answers a
reduction needs. An actual worst-case-to-average-case reduction must prove
how arbitrary successful fibers can be used, not just show that some hard
fiber exists. The nonadaptive barriers in \eref{30006686}{2} make this demand
substantive; they do not rule out every possible adaptive or non-black-box
construction.

\medskip
\textbf{Lemma 33.3 (a randomized-hardness dependency).}
If $\mathsf{NP}\subseteq\mathsf{BPP}$, then no one-way function family in
the sense of UnsolvedMath problem~3403 exists.

\emph{Proof.} Fix any deterministic polynomial-time family $f_n$ with
polynomial output length. The language
\[
 Q_f=\{\langle1^n,y,u\rangle:
       \exists x\in\{0,1\}^n\text{ extending }u,\ f_n(x)=y\}
\]
is in NP. Under the hypothesis it has a uniform randomized polynomial-time
decider. Amplify its error to at most $1/(3(n+1))$ per query. For an image
point $y$, start with the empty prefix, ask whether a preimage extending
$u0$ exists, and append zero if yes and one otherwise. With fresh coins on
each query, the conditional error bound holds even though the prefixes
are adaptive. A union bound implies all $n$ decisions are correct with
probability at least $2/3$. Then the resulting $x$ is a valid preimage.
Verify it and otherwise output failure. This is a polynomial-time inverter
with constant success on every image point, contradicting one-wayness.
\hfill$\square$

The original target therefore entails the additional implication
\begin{equation}\label{r33:dependency}
 \mathsf{NP}\subseteq\mathsf{BPP}\quad\Longrightarrow\quad
                      \mathsf P=\mathsf{NP}.
\end{equation}
Indeed, its contrapositive premise would produce an OWF, contradicting
Lemma~33.3. The elementary observation does not establish
\eqref{r33:dependency}; it warns that a merely randomized solver for an
NP-hard language is not automatically a contradiction to
$\mathsf P\ne\mathsf{NP}$. Any claimed deterministic worst-case foundation
has to address that quantifier issue.

\paragraph{Stress test.}
The hard-fiber reduction is verified only for \emph{perfect worst-case}
inversion. Replacing that oracle by an algorithm with small average success
is exactly the unjustified step being tested. The output padding keeps
image lengths and malformed outputs explicit; no recovery of the originally
sampled witness is required. Product success in \eqref{r33:smallproduct}
uses independent inputs, but the constructed inverter's validity does not
assume an arbitrary adversary treats blocks independently.

The accompanying exact tests use small truth tables for $V$, enumerate
preimages of the encoded function, and verify the rare-branch probability
and product inequalities. The finite experiments do not establish SAT
hardness; that follows from the stated reduction with a universal verifier.
The supplied source's broad reduction model is preserved, and its existing
bibliography remains unchanged. The 2010 erratum is recorded rather than
silently citing the original adaptive theorem as if its proof were intact.

\Needspace{8\baselineskip}
\paragraph{Outcome.}
\textbf{Outcome: Exploratory approach with unresolved gap.}

An explicit uniformly computable family has NP-hard worst-case inversion
but an easy uniform-input inverter, and the proposed independent-product
repair is defeated for every product length by a parameter-sensitive
inverter. The attempt also proves a randomized-hardness consequence any
positive resolution must respect. These are obstructions to particular
routes, not a counterexample to the unrestricted implication
$\mathsf P\ne\mathsf{NP}\Rightarrow\mathrm{OWF}$. No construction overcoming
the distributional exposure and deterministic-versus-randomized gaps is
supplied, and historical novelty of these elementary deductions is not
asserted.

% END RESEARCH 33

\subsection{Sources}
\begin{itemize}\raggedright
\item[\textnormal{[S1]}]\hypertarget{source:30006686:1}{} A. Akavia, O. Goldreich, S. Goldwasser and D. Moshkovitz, 'On basing one-way functions on NP-hardness', Proc. 38th Annual ACM Symp. on Theory of Computing (STOC), 2006, pp. 701-710. DOI:10.1145/1132516.1132614.
\url{https://api.openalex.org/works/doi:10.1145/1132516.1132614}.
{\small\textit{Catalogue formulation source.}}
\item[\textnormal{[S2]}]\hypertarget{source:30006686:2}{} Non-Trivial Zero-Knowledge Implies One-Way Functions.
\url{https://arxiv.org/abs/2602.17651}.
{\small\textit{Catalogue research/status reference; not a proof endorsement.}}
\item[\textnormal{[S3]}]\hypertarget{source:30006686:3}{} Hardness Along the Boundary: Towards One-Way Functions from the Worst-case Hardness of Time-Bounded Kolmogorov Complexity.
\url{https://eccc.weizmann.ac.il/report/2025/125/}.
{\small\textit{Catalogue research/status reference; not a proof endorsement.}}
\item[\textnormal{[S4]}]\hypertarget{source:30006686:4}{} Capturing One-Way Functions via NP-Hardness of Meta-Complexity.
\url{https://eccc.weizmann.ac.il/report/2023/037/}.
{\small\textit{Catalogue research/status reference; not a proof endorsement.}}
% BEGIN ADDED REFERENCES 33
\item[\textnormal{[E1]}]\hypertarget{extra:30006686:1}{} Adi Akavia, Oded
Goldreich, Shafi Goldwasser and Dana Moshkovitz, \emph{Erratum for: On
basing one-way functions on NP-hardness}, STOC 2010, p.~795.
\url{https://doi.org/10.1145/1806689.1806798}.
{\small\textit{Erratum metadata and its affected adaptive result checked
via the institutional record and [E2], 27 September 2026.}}
\item[\textnormal{[E2]}]\hypertarget{extra:30006686:2}{} Andrej Bogdanov and
Christina Brzuska, \emph{On Basing Size-Verifiable One-Way Functions on
NP-Hardness}, ECCC TR14-108 (2014), Theorem~1 and the footnote on the
unaffected nonadaptive theorem; TCC 2015, pp.~1--6.
\url{https://eccc.weizmann.ac.il/report/2014/108/download}.
\url{https://doi.org/10.1007/978-3-662-46494-6_1}.
{\small\textit{Use is of Theorem~1, not the inconsistent inclusion written
in the report abstract/corollary. Checked 27 September 2026.}}
\item[\textnormal{[E3]}]\hypertarget{extra:30006686:3}{} Suvradip Chakraborty,
James Hulett, Dakshita Khurana and Kabir Tomer, \emph{Non-Trivial
Zero-Knowledge Implies One-Way Functions}, arXiv:2602.17651v1,
19 February 2026, abstract and assumption package.
\url{https://arxiv.org/abs/2602.17651v1}.
{\small\textit{Version-specific supplement to [S2], checked
27 September 2026; no full proof audit claimed.}}

% END ADDED REFERENCES 33
\end{itemize}

% END PRESERVED SECTION
\end{document}
